\section{Discussion}
\label{sec:discussion}

\subsection{What a residual relative shift establishes}
The ESPRESSO free-shift extension improves its specified objective, and several controls retain part of that preference. A low absolute quadratic statistic does not erase the matched-model comparison. However, one gas decomposition, fixed atomic inputs and an approximate covariance do not close a systematic-error budget. The result establishes a model discrepancy without identifying a physical-frequency change.

Changing the estimator changes the conditional question. Native exposures and their combined spectrum agree under the same fixed-template estimator; shared photons and template prevent independent confirmation. Removing strong cores changes both possible model inadequacy and displacement information. The local projection quantifies the latter for one direction, without establishing what caused the observed reduction.

The model-construction concerns examined by \citet{Lee2023Models} remain relevant. Opposite-hypothesis starts and saved endpoints expose local minima but do not sample all component structures. Our merged architecture also changes bound centres. Profile searches that find lower objectives are therefore substantive revisions, not merely uncertainty calculations around an established global minimum. Small endpoint changes or successful stopping flags do not certify stationarity.

\subsection{What overlapping orders add}
An actual rest frequency cannot depend on which order records the transition. The retained 2600 contrast therefore directs attention to the observation and its model. Different sampling, masks and weights can interact with a shared imperfect gas template, so the contrast alone cannot isolate a wavelength solution, extraction step or LSF error. Matching coverage addresses one part of this problem.

Non-Gaussian ESPRESSO response is independently established \citep{Schmidt2024LSF}. Our positive kernel changes shape without changing its centroid, yet retains a contrast near $-59\,\mathrm{m\,s^{-1}}$. This rules out removal by the particular fitted family, not by all conventional responses. Sharing a kernel across traces and epochs is especially restrictive. Neither the fit nor its limited transfer performance supplies independently calibrated kernels for the historical science exposures.

The later exposures contain disjoint photons but share the gas template and reduction procedure. Barycentric velocity, detector position and observing period are also correlated. Absence of a simple univariate trend does not exclude nonlinear effects, and a trend among correlated predictors would not establish causation. Exposure-level forward modelling with external calibration addresses this remaining gap.

\subsection{From more spectra to an identifiable cross-epoch test}
Catalogue coverage, a detected multiplet and a satisfactory local fit are different readiness levels. Boundaries, truncated structure and initialization dependence cannot be discarded to obtain one regression point per absorber. Moreover, an $H_0$-adequacy filter could exclude genuinely shifted systems; the filtered diagnostic sample is not an unbiased population test of constant frequencies.

Keeping the 2382/2374 pair fixes the observable across redshifts. Unrestricted target-specific differential offsets nevertheless span its age response. This algebraic result does not imply arbitrary physical systematic amplitudes: the close pair suppresses smooth long-range distortions, and the finite slope scenarios in Sec.~\ref{sec:identifiability} produce much smaller offsets than a widely separated pair would. Independently justified bounds restore conditional information. Near-collinearity alone cannot show that modest distortions explain an arbitrarily large effect.

A shared laboratory-ratio error contributes an intercept, whereas target-specific differential terms need not. Even after controlling them, detecting a trend would not select its exponent. At these six ages, the normalized $p=1,2,3$ schedules are nearly identical after adjusting intercept and amplitude. Their small shape distances are fixed-design calculations, not measured trend evidence or physical amplitude bounds.

The inverse-logarithmic family is motivated by the author's numerical non-autonomous-map study associated with Riemann-zero data \citep{Wang2026Spectral}. It supplies no atomic Hamiltonian or Fe\,II response. An added uncertainty \emph{standard deviation} is not a deterministic level displacement: stochastic models require correlation and line-shape calculations, while drift models require level-dependent coefficients. The present diagnostics test prerequisites without asserting that physical connection.

\subsection{A prospective prediction test}
\label{sec:prospective}
A prospective extension would freeze an age-dependent prediction before examining independent validation sightlines. Keeping the same Fe\,II 2382/2374 descriptor, illustrative groups at $z=1,1.5,2$ sample cosmic ages of approximately $5.85$, $4.26$ and $3.27\,\mathrm{Gyr}$ in the adopted coordinate. Under the fixed inverse-log-square schedule, development anchors at the low and high redshifts predict
\begin{equation}
 D(1.5)_{\rm pred}=0.45756\,D_{\rm dev}(1)
                   +0.54244\,D_{\rm dev}(2).
 \label{eq:prospective_prediction}
\end{equation}
For an illustrative endpoint difference of $100\,\mathrm{m\,s^{-1}}$, the intermediate difference is $54.24\,\mathrm{m\,s^{-1}}$; neither amplitude is measured or derived from an atomic theory here. New endpoint sightlines must independently reproduce a nonzero contrast, because a constant also satisfies Eq.~\ref{eq:prospective_prediction}. Instrument, observing-condition and near-equal-redshift controls would test alternative explanations. A sufficiently precise disagreement would reject the frozen nonzero forecast; inconclusive precision would not validate it. Such a deterministic mean response is an additional hypothesis beyond an uncertainty floor. Appendix~\ref{app:prediction} specifies sample separation, covariance, failure rules, conditional power and current instrument constraints. This is a proposed protocol, with no registered prediction, selected observing programme or new observations. It does not resolve the near-degeneracy among inverse-logarithmic exponents.

\section{Conclusions}
\label{sec:conclusions}
The public observations support reproducible sensitivity tests: matched gas-model comparisons, continuum covariance controls, fixed-mask information calculations, native-order contrasts and additional archive profiles. They expose conditional inconsistencies and numerical/model dependence, while separating available spectra from calibrated frequency measurements.

A confirmation programme needs independent differential calibration, broader gas/covariance assessment, a common transition pair and selection fixed before examining shifts. Inadequate fits must remain visible. With a specified physical response, such data could constrain a time-dependent amplitude or test whether a reproduced residual follows an age law. The current study establishes neither a calibrated astrophysical frequency change nor a preferred aging exponent.
